%% Chapter 9 Solutions
%% ============================================================

\section{Chapter 9: Global Sensitivity Analysis}

%% ----------------------------------------------------------
\begin{solution}{9.1}{L3}

\solanswer

Given: $\SI{\tau_E}{R_0} = 0.6$, $\SI{\tau_I}{R_0} = 1.2$,
$\SI{\beta}{R_0} = 1.0$, $\SI{k}{R_0} = 1.0$, nominal $R_0^* = 2.0$.
Uncertainties: $\sigma_{\tau_E} = 2$ d, $\sigma_{\tau_I} = 3$ d,
$\sigma_\beta = 0.01$, $\sigma_k = 0.8$.
Nominal values: $\tau_E^* = 5.5$ d, $\tau_I^* = 7$ d, $\beta^* = 0.06$,
$k^* = 5$.

\solpart{a} \textbf{Output standard deviation (delta method).}

$\sigma_{R_0}^2 = \sum_j \left(\frac{\partial R_0}{\partial p_j}\right)^2 \sigma_{p_j}^2
= \sum_j \left(\frac{R_0^*}{p_j^*}\,\SI{p_j}{R_0}\right)^2 \sigma_{p_j}^2$.

\begin{align*}
\left(\frac{R_0^*}{\tau_E^*}\,\SI{\tau_E}{R_0}\right)^2\sigma_{\tau_E}^2
&= \left(\frac{2}{5.5}\times0.6\right)^2\times4 = (0.218)^2\times4 = 0.190\\
\left(\frac{R_0^*}{\tau_I^*}\,\SI{\tau_I}{R_0}\right)^2\sigma_{\tau_I}^2
&= \left(\frac{2}{7}\times1.2\right)^2\times9 = (0.343)^2\times9 = 1.059\\
\left(\frac{R_0^*}{\beta^*}\,\SI{\beta}{R_0}\right)^2\sigma_\beta^2
&= \left(\frac{2}{0.06}\times1.0\right)^2\times(0.01)^2 = (33.3)^2\times10^{-4} = 0.111\\
\left(\frac{R_0^*}{k^*}\,\SI{k}{R_0}\right)^2\sigma_k^2
&= \left(\frac{2}{5}\times1.0\right)^2\times(0.8)^2 = (0.4)^2\times0.64 = 0.102
\end{align*}
$\sigma_{R_0}^2 = 0.190 + 1.059 + 0.111 + 0.102 = 1.462$,
$\sigma_{R_0} \approx \mathbf{1.21}$.

\solpart{b} \textbf{Sigma-normalized indices.}

$\tilde{S}_{p_j} = (R_0^*/p_j^*)\,|\SI{p_j}{R_0}|\,\sigma_{p_j}/\sigma_{R_0}$:

\begin{center}
\begin{tabular}{lcc}
\toprule
Parameter & $\tilde{S}$ (sigma-norm.) & Contribution (\%) \\
\midrule
$\tau_I$ & $(2/7)(1.2)(3)/1.21 = 0.851$ & $58\%$ \\
$\tau_E$ & $(2/5.5)(0.6)(2)/1.21 = 0.360$ & $13\%$ \\
$\beta$  & $(2/0.06)(1.0)(0.01)/1.21 = 0.276$ & $8\%$ \\
$k$      & $(2/5)(1.0)(0.8)/1.21 = 0.264$ & $7\%$ \\
\bottomrule
\end{tabular}
\end{center}

\solpart{c} \textbf{Tornado plot.} Sorted by $\tilde{S}$: $\tau_I > \tau_E > \beta \approx k$.

\solpart{d} \textbf{Ranking comparison.}
Local SI ranking: $\tau_I$ and $k$ tied at 1.2 and 1.0; $\beta$ at 1.0; $\tau_E$ at 0.6.
Sigma-normalized ranking: $\tau_I \gg \tau_E > \beta \approx k$.

$\tau_I$ remains dominant in both, but its lead is much larger when
uncertainty is accounted for ($\sigma_{\tau_I} = 3$ d is large).
$k$ drops from second to last: despite SI $= 1.0$, its uncertainty
$\sigma_k = 0.8$ relative to $k^* = 5$ gives a coefficient of variation
of only $16\%$, while $\tau_E$ has $\sigma_{\tau_E}/\tau_E^* = 36\%$.
\textbf{$k$ changed rank most dramatically} because it has a large SI
but small relative uncertainty.

\solnote{The sigma-normalized index is the key concept here: it combines
the sensitivity (SI) with the actual uncertainty range to produce a
practical priority ranking. Students who compute only SIs without
accounting for uncertainty miss the point of this exercise.}

\end{solution}

%% ----------------------------------------------------------
\begin{solution}{9.2}{L3}

\solanswer

\solpart{a} \textbf{LHS implementation and comparison.}

\begin{verbatim}
rng(42);
N = 100; m = 3;
X_lhs = lhs_sample(N, m, zeros(1,m), ones(1,m));
X_rnd = rand(N, m);
\end{verbatim}

\solpart{b} \textbf{Scatter plots.} The LHS scatter plots show more
uniform coverage with fewer clusters; the random sample shows visible gaps
and clusters in all three pairwise projections.

\solpart{c} \textbf{Mean and standard deviation.}
Both LHS and random sampling should give means near $0.5$ and
standard deviations near $1/\sqrt{12} \approx 0.289$ for a uniform
distribution on $[0,1]$. LHS is typically closer to both exact values
because the stratification ensures coverage of the full range. For $N = 100$,
the improvement in mean accuracy is $\sim$30--50\% and in standard deviation
$\sim$20--40\%.

\solpart{d} \textbf{Histogram comparison.}
The LHS histogram shows a much flatter, more uniform distribution
(each of the 100 strata has exactly one sample). The random sample
histogram shows the typical Poisson variation (some bins have 0 or 3+ samples).
For sensitivity analysis, the LHS coverage is preferable because
it ensures the model is evaluated in every region of the parameter space.

\solnote{Part (c) is the quantitative check: students should compute the
actual mean and std and report how close they are to the theoretical values.
The LHS improvement in mean accuracy comes from the fact that the sample
mean of a stratified sample is an unbiased estimator with lower variance
than the sample mean of a random sample (stratification theorem).}

\end{solution}

%% ----------------------------------------------------------
\begin{solution}{9.3}{L3}

\solanswer

\solpart{a-b} \textbf{LHS + PRCC computation.}

Using $N = 200$, $\pm50\%$ around nominal values:
\begin{verbatim}
p = sir_nominal();
lb = 0.5*[p.k, p.beta, p.tau, p.L];
ub = 1.5*[p.k, p.beta, p.tau, p.L];
rng(42); X = lhs_sample(200, 4, lb, ub);
Y = arrayfun(@(i) sir_I_max(X(i,:), p), 1:200)';
[rho, pval] = compute_prcc(X, Y, {'k','beta','tau','L'});
\end{verbatim}

Expected PRCC values (approximate):
\begin{center}
\begin{tabular}{lcc}
\toprule
Parameter & PRCC & $p$-value \\
\midrule
$k$     & $+0.85$ & $< 0.001$ \\
$\beta$ & $+0.84$ & $< 0.001$ \\
$\tau$  & $+0.72$ & $< 0.001$ \\
$L$     & $-0.04$ & $> 0.5$  \\
\bottomrule
\end{tabular}
\end{center}

\solpart{c} \textbf{PRCC bar chart.} Sorted: $k \approx \beta > \tau \gg L$ (similar to SI ranking).

\solpart{d} \textbf{Strongest negative PRCC: $L$.}
$L$ (mean lifespan) has a negative PRCC: longer lifespans reduce peak infections.
In plain language: populations with longer lifespans have smaller background
mortality rates, which means the susceptible pool is replenished more slowly,
and a larger fraction of the population is already immune/recovered by the
time of peak infections. The effect is very small ($|\text{PRCC}| \approx 0.04$,
not statistically significant).

\solpart{e} \textbf{Comparison with local SI tornado plot.}
Both methods rank $k$ and $\beta$ highest, $\tau$ third, $L$ negligible.
The PRCC ranking agrees with the local SI ranking for this model.
Agreement is expected because peak infections $I_{\max}$ is monotone in
each parameter over the $\pm50\%$ uncertainty range (the model is
not near any threshold), satisfying the PRCC assumption.

\solnote{For the SIR model with $R_0 = 2.1$ and $\pm50\%$ variation,
the PRCC and local SI rankings should agree well. Discrepancies would
appear if: (a) $\beta$ variation includes $R_0 < 1$ (non-monotone at threshold),
or (b) parameter interactions are strong (nonlinear effects). Students
should note that agreement here is specific to this model and range.}

\end{solution}

%% ----------------------------------------------------------
\begin{solution}{9.4}{L4}

\solanswer

$Y = 2Q_1^2 + Q_2$, $Q_1, Q_2 \sim \mathcal{U}(0,1)$, independent.

\solpart{a} \textbf{Variance decomposition.}
$\E[Y] = 2\E[Q_1^2] + \E[Q_2] = 2/3 + 1/2 = 7/6$.
$\Var(Y) = 4\Var(Q_1^2) + \Var(Q_2)$.
$\Var(Q_1^2) = \E[Q_1^4] - (\E[Q_1^2])^2 = 1/5 - 1/9 = 4/45$.
$\Var(Q_2) = 1/12$.
$\Var(Y) = 4(4/45) + 1/12 = 16/45 + 1/12 = 64/180 + 15/180 = 79/180$.

$D_1 = \Var[\E(Y|Q_1)] = \Var[2Q_1^2 + 1/2] = 4\Var(Q_1^2) = 16/45$.
$D_2 = \Var[\E(Y|Q_2)] = \Var[2/3 + Q_2] = 1/12$.

\solpart{b} \textbf{Sobol indices.}
$S_1 = D_1/\Var(Y) = (16/45)/(79/180) = (16/45)(180/79) = 64/79 \approx 0.810$.
$S_2 = D_2/\Var(Y) = (1/12)/(79/180) = (180/12)/79 = 15/79 \approx 0.190$.
$S_1 + S_2 = 79/79 = 1.0$. Yes, they sum to 1, because the model is
\emph{additive} (no interaction terms), so the full variance is explained
by the first-order indices.

\solpart{c} \textbf{Total indices.}
For an additive model: $S_{T_1} = S_1 = 64/79$, $S_{T_2} = S_2 = 15/79$.
Total indices equal first-order indices when there are no interactions.

\solpart{d} \textbf{Adding interaction $Q_1 Q_2$: $Y' = 2Q_1^2 + Q_2 + Q_1 Q_2$.}
Without computing: adding $Q_1 Q_2$ introduces an interaction term.
$S_1$ decreases (some variance is now shared with $Q_2$); $S_2$ decreases
for the same reason. The total indices $S_{T_1} > S_1$ and $S_{T_2} > S_2$
because each parameter contributes to the interaction term.

\textit{Verification:}
$\Var(Y') = 4\Var(Q_1^2) + \Var(Q_2) + \Var(Q_1 Q_2) + 2\text{Cov}(\ldots)$
(with $\text{Cov}$ terms vanishing by independence).
$\Var(Q_1 Q_2) = \E[Q_1^2 Q_2^2] - (\E[Q_1 Q_2])^2 = (1/3)(1/3) - (1/2)^2(1/4)
= 1/9 - 1/16 = 7/144$.
$D_{12} = \Var(Y') - D_1' - D_2' - $ (remaining interaction variance).
Computing exactly: $S_1' < 0.810$, $S_2' < 0.190$, $S_1' + S_2' < 1$.
The interaction term $D_{12} > 0$ accounts for the remainder.

\solnote{Part (b)'s exact sum-to-1 for the additive model is a crucial
result: for any purely additive model, all first-order indices sum to 1
and total indices equal first-order indices. Part (d) demonstrates the
effect of interactions without requiring exact computation, testing
conceptual understanding over algebraic facility.}

\end{solution}

%% ----------------------------------------------------------
\begin{solution}{9.5}{L4}

\solanswer

Using \texttt{run\_global\_sa\_sir} from the repository with $N = 2000$.

\solpart{a} \textbf{Sobol indices (nominal parameters, $\pm50\%$ range).}

Approximate values from Jansen estimators:
\begin{center}
\begin{tabular}{lcccc}
\toprule
Param & $\hat{S}_i$ & $\hat{S}_{T_i}$ & $\hat{S}_{T_i} - \hat{S}_i$ \\
\midrule
$k$     & $0.38$ & $0.43$ & $0.05$ \\
$\beta$ & $0.38$ & $0.43$ & $0.05$ \\
$\tau$  & $0.20$ & $0.25$ & $0.05$ \\
$L$     & $0.00$ & $0.00$ & $0.00$ \\
\bottomrule
\end{tabular}
\end{center}

\solpart{b} \textbf{Parameters with $S_{T_i} \gg S_i$.}
For all three active parameters, $S_{T_i} - S_i \approx 0.05$, indicating
moderate interaction effects (primarily $k$-$\beta$ and $k$-$\tau$ interactions).
No parameter has $S_{T_i} \gg S_i$ for this model at this range;
interactions are present but modest. This is consistent with the SIR model
being a smooth, moderately nonlinear ODE.

\solpart{c} \textbf{Sums.}
$\sum S_i \approx 0.96 < 1$ (interaction terms absorb the remaining $\sim4\%$).
$\sum S_{T_i} \approx 1.11 > 1$ (interactions are double-counted in total indices).
For a purely additive model, both sums would equal 1. The observed values
confirm the presence of small but nonzero interactions.

\solpart{d} \textbf{Best parameter to measure precisely.}
To reduce $\Var(I_{\max})$ most, measure the parameter with the largest
$S_{T_i}$ (total index), since $S_{T_i}$ measures the variance that
would remain if all uncertainty in parameter $i$ were eliminated, including
interactions. Here $k$ and $\beta$ are tied at $S_{T_i} \approx 0.43$.
Measuring either $k$ or $\beta$ precisely reduces $\Var(I_{\max})$
by $\approx 43\%$ (more precisely, the variance conditional on knowing $k$ exactly
is $(1 - S_{T_k})\Var(I_{\max}) \approx 0.57\,\Var(I_{\max})$).

\solnote{Part (d) is a key policy application of total indices: the total
index, not the first-order index, answers ``which parameter measurement
would most reduce output uncertainty?'' First-order indices answer
``which parameter contributes most to variance independently?'' Both
are important, and students should be able to distinguish them.}

\end{solution}

%% ----------------------------------------------------------
\begin{solution}{9.6}{L4}

\solanswer

Using $r = 15$ Morris trajectories, $p = 4$ levels ($\Delta = 4/6 = 2/3$),
$\pm50\%$ parameter ranges.

\solpart{a} \textbf{Morris measures.} Approximate values from \texttt{morris\_screening}:
\begin{center}
\begin{tabular}{lccc}
\toprule
Param & $\mu^*$ & $\sigma$ & Classification \\
\midrule
$k$     & $\sim 80$ & $\sim 20$ & Important, nonlinear \\
$\beta$ & $\sim 78$ & $\sim 19$ & Important, nonlinear \\
$\tau$  & $\sim 42$ & $\sim 14$ & Important, nonlinear \\
$L$     & $\sim 1$  & $\sim 0.5$& Negligible \\
\bottomrule
\end{tabular}
\end{center}

\solpart{b} \textbf{Diagnostic plot.}
The $\mu^*$ vs.\ $\sigma$ plot shows $k$ and $\beta$ in the upper-right
(large $\mu^*$ and $\sigma$: important and nonlinear), $\tau$ in the
middle-right, and $L$ near the origin (negligible).

\solpart{c} \textbf{Most and least influential.}
Most: $k$ (or $\beta$, tied). Least: $L$.

\solpart{d} \textbf{Reduced Sobol (3 parameters, $L$ fixed).}

With $L$ fixed: Sobol cost $= N(m+2) = 2000\times5 = 10{,}000$ evaluations
vs.\ full $2000\times6 = 12{,}000$. Cost reduction: 17\%.
Approximate results: $S_k \approx 0.43$, $S_\beta \approx 0.43$, $S_\tau \approx 0.24$;
total indices similar. Fixing $L$ does not materially change the other
indices (since $L$ was negligible). The Morris screening correctly identified
that $L$ could be dropped.

\solnote{The Morris cost here is $15\times(4+1) = 75$ evaluations
vs.\ $12{,}000$ for full Sobol — a $160\times$ cost reduction for the
screening phase. This illustrates Morris's value as a cheap first step
that identifies which parameters can be dropped from the expensive Sobol analysis.}

\end{solution}

%% ----------------------------------------------------------
\begin{solution}{9.7}{L5}

\solanswer

\textbf{This exercise requires the student to choose and describe a specific
published model.} The model answer below uses the Chitnis dengue model~\cite{chitnis2008determining}
as an example; any 8+ parameter epidemic model with published parameter estimates is acceptable.

\solpart{a} \textbf{Model, parameters, uncertainty ranges.}
Chitnis et al.\ (2008) dengue model: 8 state variables, 14 parameters including
mosquito biting rate $b_v$, transmission probabilities $\beta_{vh}$ and $\beta_{hv}$,
mosquito and human mortality rates, and mean durations. Uncertainty ranges:
$\pm50\%$ for biological parameters (based on reported ranges in Table 2 of the paper);
$\pm30\%$ for demographic parameters (population data are more precisely known).

\solpart{b} \textbf{QOI justification.}
Primary: $R_0$ (determines whether dengue can persist in the population —
the most policy-relevant threshold). Secondary: equilibrium human infection
prevalence $\bar{I}_h/N_h$ (determines endemic burden for healthcare planning).

\solpart{c} \textbf{PRCC vs.\ Sobol choice.}
Use \textbf{Sobol indices}. The dengue model is moderately nonlinear and
has known interaction between biting rate $b_v$ and transmission probabilities
($R_0 \propto b_v^2\beta_{vh}\beta_{hv}$, so interactions are structural).
PRCC assumes monotone relationships; near threshold or at high $R_0$,
this assumption may not hold for all parameters. Sobol indices capture
interactions and do not require monotonicity.

\solpart{d} \textbf{Sample size $N$.}
Use $N = 2000$ for each Sobol matrix. Run the convergence check:
compute Sobol indices at $N = 200, 500, 1000, 2000$; report the
value of $N$ at which the first-order indices change by $< 2\%$ as $N$ doubles.
For a 14-parameter model, the total evaluation cost is $N(14+2) = 32{,}000$ runs.
If each run takes 2 minutes, this is $\approx 4.4$ days of computation on one core.

\solpart{e} \textbf{Presentation format.}
Two tornado plots of Sobol $S_{T_i}$ (one per QOI), with parameters
labeled as interventions: ``Reducing mosquito biting rate by $X\%$ through
bed nets reduces $R_0$ by approximately $Y\%$.'' A two-panel figure
showing which parameters appear in both tornado plots (co-optimization)
and which are QOI-specific. Use traffic-light color coding: red for
high-priority, green for low-priority intervention targets.

\solpart{f} \textbf{Potential misleading results.}
If the model's best-fit parameters give $R_0$ near 1 (endemic/non-endemic
threshold), Sobol indices at those nominal values will be dominated by
parameters near the threshold. A 2\% uncertainty in $\beta_{vh}$ might
shift the model from endemic to non-endemic, giving extremely large apparent
sensitivity for $\beta_{vh}$. Recommendation: compute Sobol indices at
three nominal parameter sets ($R_0 = 1.5$, $2.0$, $3.0$) and report
whether the parameter rankings are stable across this range.

\solnote{This is an open-ended design exercise. Full credit requires
all six components with operationally specific content. Vague statements
(``choose the right sample size,'' ``present clear figures'') without
specifics receive partial credit. The model choice affects all answers;
a reasonable choice with internally consistent answers earns full credit
even if different from the example above.}

\end{solution}

%% ----------------------------------------------------------
\begin{solution}{9.8}{L6}

\solanswer

\textbf{Methodological concerns with Sobol analysis ignoring correlation
among three parameters.}

\solpart{a} \textbf{Mathematical problem.}

Standard Sobol indices assume parameter \emph{independence}: the matrices
$A$ and $B$ are sampled independently from the joint distribution $p(\vec{p})
= \prod_j p_j(p_j)$. When three parameters are correlated, the correct
joint distribution is $p(\vec{p}) = p(p_1, p_2, p_3) \times \prod_{j>3} p_j(p_j)$,
which cannot be factored. Sampling independently from marginal distributions
and then compositing (the Saltelli method) generates non-physical parameter
combinations (e.g., high phosphorus loading but low nitrogen loading, which
does not occur in the real watershed system). The resulting Sobol indices
estimate sensitivity to artificially independent variations, not to the
actual coupled variations that the system experiences.

\solpart{b} \textbf{Effect on each correlated parameter's index.}

\textit{Phosphorus loading $P$:} The standard $S_1(P)$ likely \textit{overestimates}
the true first-order index. When $P$ and $N$ are positively correlated,
the actual variation of $P$ is accompanied by variation in $N$. The standard
estimator attributes only the $P$ variation to $S_1(P)$; the interaction
that would occur under the true distribution is ignored, artificially
inflating the apparent individual contribution of $P$.

\textit{Nitrogen loading $N$:} Same logic applies: $S_1(N)$ likely overestimates.

\textit{Water residence time $T_w$:} If $T_w$ is correlated with both $P$ and $N$
(longer residence = more loading), $S_1(T_w)$ may be \textit{underestimated}
because the positive correlation means that varying $T_w$ while holding $P$
and $N$ fixed understates the total effect of residence time in the real system.

\solpart{c} \textbf{Minimal correction.}

Use a copula-based sampling method that preserves the observed correlations.
The Iman--Conover rank correlation constraint method is the simplest:
(1) generate independent LHS samples; (2) apply the Iman--Conover algorithm
to induce the observed rank correlations among the three parameters;
(3) compute Sobol indices using the correlated sample matrices.
This adds one matrix transformation step to the existing workflow with
no additional model evaluations.

\solpart{d} \textbf{Why sum-to-1 is not a valid independence check.}

The Sobol first-order indices sum to 1 only for an additive model
with independent parameters. With correlated parameters, the indices
can still sum to approximately 1 even with incorrect independence assumptions,
because the sum reflects the total explained variance, which equals 1 by
construction (the variance is fully accounted for by some combination of
first-order and interaction terms). The sum being 1 says nothing about
whether the individual index attributions are correct — it is a property
of the decomposition algorithm, not of the physical model or the sampling design.

\solnote{This is a rigorous L6 exercise testing critical evaluation of
a real methodological error. Part (d) is the subtlest: students often
think ``indices sum to 1'' is a validity check, but it is not — it is
simply a consequence of the ANOVA decomposition. Full credit requires
a clear mathematical explanation, not just ``the check doesn't work.''}

\end{solution}
